\subsection{Da Geometria e Movimento}
\par §1 --- Dado a concepção e compreensão das quantidades e relações presentes em um sistema, é necessário
aplicá-las em suas devidas funções e para implementação formal das mesmas. 
Propriedades como \textbf{tamanho}, \textbf{posição} e \textbf{variação} tomam forma e se 
tornam características fundamentais deste sistema.

\begin{figure}[h]
    \centering
    \includegraphics[width=8cm]{Martelo2}
\end{figure}

\par
§ --- Sucedendo o exemplo anterior, o martelo e a corda, anteriormente \textit{meros conceitos}, tomam \textbf{espaço} e 
\textbf{forma} em um campo cartesiano bidimensional. De mesma forma, é implicado a 
existência de movimento \\$S = S_1 - S_0$ em um período de tempo $\Delta{t} = t_1 - t_0$.
